\begin{textsample}{POS Dim 4 – global\_north\_2023\_10 – Score 12.00 – global\_north\_2023\_10/91686283.txt}  \label{ex:f4_pos_450}
Israeli \textbf{Prime} Minister Benjamin Netanyahu addresses the 78th United Nations General Assembly at the UN headquarters in New York on September 22, 2023 File : Reuters/Mike Segar Israeli \textbf{Prime} Minister Benjamin Netanyahu has spent the past three decades sounding the alarm about Iran ' s nuclear programme and threatening to attack the country on countless occasions.
Most recently in September, he said in a speech to the United Nations General Assembly that Tehran must face a"credible nuclear threat"before his office corrected the record to"credible military threat".
After Hamas ' s attack on October 7, Netanyahu may finally be able to act on his threats.
The gruesome scenes in southern Israel have provided the Israeli \textbf{prime} minister with the necessary pretext and international backing for a wider response.
Netanyahu has both a political and a personal stake in all this.
A drawn-out regional conflict would block or at least postpone any official accountability for his utter failure to prevent Hamas ' s attack from happening in the first place and could also put his multiple indictments on corruption charges on an indefinite hold. @ @ @ @ @ @ @ @ @ @ \textbf{prime} minister to a wartime leader, with opposition parties clamouring to join him in a national unity government.
He declared war and ordered an immediate retaliation against Hamas ' s stronghold in Gaza.
The Israeli army unleashed a vicious campaign of bombardment on the overpopulated Gaza Strip, killing more than 500 people, and preparing for a potential land invasion.
Netanyahu has not elaborated on the next \textbf{phases} of the war, but he has received the unconditional support of Western governments to do what it takes, as long as it takes, to"defend Israel".
The administration of US \textbf{President} Joe Biden has gone even further, providing Israel with more arms and ammunition, dispatching its most modern and sophisticated aircraft carrier, the Ford, along with a number of destroyers to the Eastern Mediterranean, and beefing up other forces stationed in the region, enough to start World War III.
Biden ' s motivation for the escalatory deployment is, reportedly, strategic deterrence, meant to ensure that"no enemies of Israel @ @ @ @ @ @ @ @ @ @.
But historically, Israel has never allowed any foreign boots on its soil, and is in no need of the US armadas to take on Hamas.
Biden ' s incentive, therefore, could also be political, ie to ensure that the GOP doesn't exploit the Israeli drama at his expense ahead of the presidential elections in 2024.
Already, Republican opponents have tried to link Biden ' s recent \textbf{prisoner} \textbf{swap} \textbf{deal} with Iran, which involved the unfreezing of \$6bn in Iranian assets, to the Hamas attacks.
But Netanyahu and his fanatic ministers may have something very different in mind for the US deployment, that goes beyond military deterrence and political posturing.
He may try to widen the scope of the war to include Iran.
His government has already accused Iran of supporting and directing Hamas ' s operation, as it has previously done about other Palestinian attacks on Israelis.
Scores of Israel supporters and neoconservatives, as well as mediapundits in the US and Europe, have joined in by making the case for @ @ @ @ @ @ @ @ @ @ -- based on interviews with unnamed local sources -- that Iranian officials and members of the Islamic Revolutionary Guards Corps were directly involved in orchestrating and \textbf{planning} the attacks over several weeks.
US officials have said they haven't seen evidence of Tehran ' s involvement, yet.
For its part, Iran has called the attack a spontaneous Palestinian action in self-defence, but officials have not tried to hide their glee at Israel ' s misfortune.
They have expressed confidence that the attack will deter further Arab, meaning Saudi, normalisation with Israel, and eventually lead to its downfall.
Iran and its allies ' temerity may well come back to haunt them, just as Israel ' s hubris did -- leading to its utter humiliation at the hands of Hamas fighters.
Neither Iran nor Israel is learning from history, as they continue to escalate their proxy conflict towards war.
For years, the Israeli army and secret services have sabotaged the Iranian nuclear programme and targeted Iranian assets abroad.
Iran for its part has supported various @ @ @ @ @ @ @ @ @ @ and Israeli allies.
Despite his bluster and bravado, Netanyahu couldn't and wouldn't attack Iran, without a green light and support from the US.
But the bloody attacks are a game-changer, giving the Israeli \textbf{prime} minister the perfect opportunity to realise his fantasy of crushing Iran, by tricking the Biden administration into war.
This will not be easy considering Biden ' s presumed commitment to end"the forever wars", reflected in the humiliating \textbf{withdrawal} from Afghanistan in 2021.
His administration has also moved to prioritise the great power competition with China and Russia, especially after the latter ' s invasion of Ukraine.
But in reality, the US has not \textbf{withdrawn} from the Middle East, it has merely moved around its forces and military assets in the region.
Biden himself has vowed to"not walk away and leave a vacuum to be filled by China, Russia, or Iran".
Once the case against Tehran ' s role in the attacks has been fully articulated by Israel and the US @ @ @ @ @ @ @ @ @ @ the \textbf{release} of Israeli \textbf{captives} taken by Hamas -- a top priority for Netanyahu.
If Iran refuses and chooses to use Hezbollah as leverage against Israel, this could well trigger a wider confrontation that draws in the US with incalculable consequences.
Unfortunately, in the adulterated world of Washington politics, unconditional US support of Israel is the only thing that Republican and Democrats \textbf{agree} on.
It is crucial to remember that the situation in 2023 is vastly more challenging and complicated than the lead-up to the invasion of Iraq in 2003, which ended in utter disaster for the US and Iraqis.
A repeat against Iran is sure to be far worse for all concerned.

% matched lemmas: agree, captive, deal, phase, plan, president, prime, prisoner, release, swap, withdraw, withdrawal
\end{textsample}
